% !TEX root = ../main-es.tex
%TC:macro \limite [ignore,ignore]
%TC:macro \nota [ignore]
%TC:macro \TODO [ignore]
%TC:macro \parencite [ignore]
%TC:macro \textcite [ignore]
\section{Marco conceptual}\label{sec:marco}

\nota{Numeración de secciones de R\&N según la 4.ª edición (2020); info/01-mapping.md pide cotejarla contra el ejemplar físico antes de entregar.}

El trabajo descansa sobre cinco cuerpos de teoría, todos presentes en el curso, que aquí se usan como piezas funcionales del modelo y no como adorno: La caracterización de agentes y entornos de tarea, la planificación en entornos que cambian (en lazo abierto o en lazo cerrado), el costo de la deliberación, los agentes basados en conocimiento y la inferencia probabilística. Cabe mencionar que la búsqueda informada, que no recibe una subsección propia, también es una pieza funcional, pues A* es el planificador que el agente invoca cada vez que construye o repara un plan.

\subsection{Agentes, entornos de tarea y el orden de la taxonomía}

\textcite[§2.3]{russellNorvig2020} especifican un entorno de tarea con la descripción PEAS (medida de rendimiento, entorno, actuadores y sensores) y lo clasifican según seis propiedades, a saber, si es total o parcialmente observable, determinista o estocástico, episódico o secuencial, estático o dinámico, discreto o continuo, y de uno o varios agentes. Sobre esa caracterización construyen una taxonomía de programas de agente de sofisticación creciente (§2.4), que va del agente reflejo simple (que responde solo a la percepción actual) al reflejo basado en modelo (que mantiene un estado interno del mundo), al basado en objetivos (que planifica una secuencia de acciones hacia una meta) y al basado en utilidad. Vale la pena anotar que la tesis del capítulo no es que el agente más sofisticado sea siempre el mejor, sino que el programa adecuado depende del entorno de tarea, lo cual es justamente lo que este trabajo pone a prueba al medir en qué región del espacio de parámetros cada arquitectura gana.

El extremo reactivo de la taxonomía tiene además autor y manifiesto propios, pues \textcite{brooks1991intelligence} sostuvo que un agente acoplado al mundo puede prescindir de representaciones y planes explícitos si sus lazos de percepción y acción son lo bastante rápidos, de forma que el mundo sirva como su propio modelo. Esa postura y la de la planificación deliberativa conviven en el experimento como políticas comparables, sin que ninguna se dé por buena de antemano.

\subsection{Planificar y actuar en un mundo que cambia}

Un plan clásico se compila sobre un modelo del mundo y se ejecuta sin realimentación, lo cual es óptimo en un entorno estático, determinista y totalmente observable, y frágil en cuanto alguna de esas condiciones falla \parencite[§11.5]{russellNorvig2020}. Para esos entornos el libro distingue la planificación sin sensores, la planificación contingente (que ramifica el plan según lo que se observe) y la planificación en línea (§11.5.3), en la que el agente monitorea la ejecución y replanifica cuando el mundo deja de comportarse como el plan supone. El monitoreo puede hacerse en tres niveles: De acción, cuando se verifica que la precondición de la siguiente acción se cumple; de plan, cuando se verifica que el resto del plan todavía puede alcanzar la meta; y de objetivo, cuando se verifica si conviene perseguir otra meta. En el lenguaje de control, el plan ejecutado sin realimentación es un \emph{lazo abierto} y el plan monitoreado y reparado durante la ejecución es un \emph{lazo cerrado}.

La tradición de planificación aportó además dos advertencias que este trabajo se propone medir en lugar de suponer. Primeramente, reusar o reparar un plan es, en el caso general, tan difícil computacionalmente como generarlo de nuevo \parencite{nebel1995plan}, por lo que la intuición de que reparar sale más barato que replanificar hay que demostrarla en cada dominio. Progresivamente, la distinción entre replanificar desde cero y reparar el mínimo necesario se formuló junto con la estabilidad del plan como métrica \parencite{fox2006plan}, y se acompañó de algoritmos que replanifican de forma incremental en retículas donde el mundo cambia mientras el agente se mueve, como D* Lite \parencite{koenig2002dlite}, o que entregan una solución aceptable de inmediato y la mejoran si queda tiempo, como ARA* \parencite{likhachev2003ara}. Finalmente, la pregunta de cuándo conviene saltarse una replanificación tiene ya tratamiento formal en procesos de decisión parcialmente observables, tanto desde el control predictivo \parencite{ma2018performance} como desde la simplificación en lazo abierto con garantías de desempeño \parencite{kong2026open}. Lo que todas estas líneas comparten es un replanificador exacto, de forma que lo que la era de los modelos de lenguaje añade no es la replanificación en sí sino la naturaleza del replanificador, que pasa a ser un sistema estocástico con una competencia que no se conoce de antemano y un costo por consulta.

\subsection{El costo de deliberar}

La tradición del metarazonamiento y la racionalidad limitada trata la deliberación como una acción con costo que compite con actuar \parencite{russell1991right,zilberstein2011metareasoning}: Pensar más mejora la decisión, pero el tiempo que se dedica a pensar es tiempo en el que el agente no actúa y, en un entorno dinámico, tiempo durante el cual la información sobre la que se piensa caduca. De ahí se siguen dos consecuencias que el modelo hace explícitas. La primera es que un plan calculado sobre el mundo tal como estaba al empezar a deliberar puede nacer desactualizado, y la segunda es el \emph{thrashing}, el modo de fallo en el que el agente reinicia una y otra vez una deliberación que una novedad del mundo interrumpe antes de terminar, de forma que paga por pensar sin llegar a actuar. Un tope fijo de iteraciones, como el de la instancia que motivó el trabajo, es una política de deliberación puesta por decreto, y la alternativa que propone esta tradición es derivarla de lo que la deliberación cuesta y de lo que compra.

La historia de la auto-modificación se puede ordenar sobre ese mismo eje. El programa de damas de \textcite{samuel1959some} ajustaba los pesos de su propia función de evaluación a partir de las partidas que jugaba, el \emph{chunking} de SOAR compila experiencia en reglas nuevas que se insertan mientras el sistema corre, con un costo marginal casi nulo \parencite{laird2012soar}, y la máquina de Gödel reescribe cualquier parte de su código solo cuando encuentra una prueba de que la reescritura mejora su utilidad, es decir, con competencia perfecta y un costo de deliberación astronómico \parencite{schmidhuber2003goedel}. Entre esos dos extremos viven todos los reparadores reales, y el plano formado por la competencia $p$ y el costo $c$ del oráculo es precisamente el espacio que los contiene.

\subsection{Agentes basados en conocimiento}

Un agente basado en conocimiento mantiene una base de conocimiento (KB) de oraciones en un lenguaje formal y opera con dos primitivas, \texttt{TELL}, que añade a la KB lo que el agente percibe, y \texttt{ASK}, que consulta qué se sigue de lo que sabe \parencite[§7.1]{russellNorvig2020}. \textcite{newell1982knowledge} propuso describir a un agente así en el \emph{nivel de conocimiento}, es decir, por lo que sabe y por sus metas, con independencia de cómo esté implementado, lo cual justifica que la misma KB pueda realizarse con un demostrador de teoremas o con estructuras de datos más eficientes, siempre que ambas respondan lo mismo a las mismas preguntas.

La consulta se resuelve por implicación lógica: Una oración $\alpha$ se sigue de la KB ($\mathit{KB} \models \alpha$) si es verdadera en todos los modelos de la KB, lo cual puede decidirse por enumeración de modelos (§7.4) o, de forma mucho más eficiente en la práctica, con el algoritmo DPLL (§7.6), que busca un modelo de $\mathit{KB} \wedge \neg\alpha$ con retroceso, propagación unitaria y símbolos puros \parencite{davis1962machine}. Cuando ni $\alpha$ ni $\neg\alpha$ se siguen de la KB, la respuesta honesta es \emph{desconocido}, que es el tercer valor que el agente de este trabajo usa para decidir si repara su plan. La implicación es además \emph{monótona} (§7.5), pues si $\mathit{KB} \models \alpha$ entonces $\mathit{KB} \wedge \beta \models \alpha$ para cualquier $\beta$, de forma que añadir conocimiento nunca invalida una conclusión.

La monotonía parece chocar con un mundo que cambia, pues lo que era cierto deja de serlo. La salida del libro es indexar los símbolos por el tiempo (\emph{fluentes}, §7.7.1), de modo que $\mathit{En}(r,c,t)$ afirma que el recurso $r$ estaba en la celda $c$ en el tick $t$, y escribir la dinámica del mundo como axiomas de estado sucesor que relacionan los fluentes de un tick con los del siguiente. Con esa representación, mantener el conjunto de estados compatibles con todo lo observado es una estimación de estado lógica, la misma que \textcite[§4.4]{russellNorvig2020} plantean para la búsqueda con observaciones parciales, y que se actualiza en dos pasos, la predicción con el modelo de transición y la poda con la observación.

\subsection{Inferencia probabilística y decisión}

Donde la lógica dice qué es posible, la probabilidad dice qué es probable. El \emph{filtrado} \parencite[§14.2]{russellNorvig2020} mantiene la distribución sobre el estado oculto dada toda la evidencia y la actualiza en cada paso con la misma estructura de predicción y corrección,
\begin{equation}
  P(X_{t+1} \mid e_{1:t+1}) \;=\; \alpha\, P(e_{t+1} \mid X_{t+1}) \sum_{x_t} P(X_{t+1} \mid x_t)\, P(x_t \mid e_{1:t}),
\end{equation}
con $\alpha$ una constante de normalización, de forma que el soporte de la distribución filtrada es exactamente el conjunto de estados que la estimación lógica no puede descartar. Cuando un parámetro del modelo de transición no se conoce, el aprendizaje bayesiano de parámetros (§20.2) lo trata como una variable más: Para una tasa de Bernoulli, una distribución previa $\mathrm{Beta}(a,b)$ es conjugada, y tras observar $m$ éxitos en $n$ ensayos la posterior es $\mathrm{Beta}(a+m,\, b+n-m)$, cuya media es $(a+m)/(a+b+n)$. Para decidir, el principio de utilidad esperada máxima (§16.1) elige la acción que maximiza $\mathit{EU}(a \mid e) = \sum_{s'} P(s' \mid a, e)\, U(s')$, con el pequeño detalle de que la decisión es tan buena como el modelo de costos con que se evalúa $U$, que es justamente donde falla el agente bayesiano de este trabajo.

Dos herramientas más, ambas de la sesión de inferencia probabilística para modelos basados en agentes, completan el marco. La computación bayesiana aproximada (ABC) por rechazo permite inferir los parámetros de un simulador cuya verosimilitud no se puede escribir pero sí simular \parencite{beaumont2002approximate}: Se muestrean parámetros de la distribución previa, se simula con cada uno, y se aceptan los que producen estadísticos resumen cercanos a los observados, de forma que las muestras aceptadas aproximan la posterior. El análisis de sensibilidad por descomposición de la varianza, por último, mide qué fracción de la varianza de una respuesta explica cada factor barrido con el cociente $\eta^2 = SS_{\text{entre}}/SS_{\text{total}}$, es decir, con la varianza que desaparece al reemplazar cada corrida por la media de su nivel.
